% Faithful transcription of official CCC2023 article15, Sections2-3 (pp.15:11-15:15).
% Original author numbering is written explicitly; no new proofs or hypotheses.
\section*{2 Preliminaries}
\subsection*{2.1 Relative Rank and its Properties}
We first describe the notation that we need to define the measures that we use to prove Theorems 1 and 3.
\paragraph{Definition 5 (Relative Rank Measure of [23,36]).} Let $f$ be a polynomial that is set-multilinear with respect to the variable partition $(X_1,X_2,\ldots,X_d)$ where each set is of size $n$. Let $w=(w_1,w_2,\ldots,w_d)$ be a tuple (or word) of non-zero real numbers such that $2^{|w_i|}\in[n]$ for all $i\in[d]$. For each $i\in[d]$, let $X_i(w)$ be the variable set obtained by removing arbitrary variables from the set $X_i$ such that $|X_i(w)|=2^{|w_i|}$, and let $X(w)$ denote the tuple of sets of variables $(X_1(w),\ldots,X_d(w))$. Corresponding to a word $w$, define $P_w:=\{i\mid w_i>0\}$ and $N_w:=\{i\mid w_i<0\}$. Let $M_w^P$ be the set of all set-multilinear monomials over a subset of the variable sets $X_1(w),X_2(w),\ldots,X_d(w)$ indexed by $P_w$, and similarly let $M_w^N$ be the set of all set-multilinear monomials over these variable sets indexed by $N_w$.
Define the `partial derivative matrix' matrix $M_w(f)$ whose rows are indexed by the elements of $M_w^P$ and columns indexed by the elements of $M_w^N$ as follows: the entry of this matrix corresponding to a row $m_1$ and a column $m_2$ is the coefficient of the monomial $m_1\cdot m_2$ in $f$. We define
\[\operatorname{relrk}_w(f):=\frac{\operatorname{rank}(M_w(f))}{\sqrt{|M_w^P|\cdot|M_w^N|}}=\frac{\operatorname{rank}(M_w(f))}{2^{\frac12\sum_{i\in[d]}|w_i|}}.\]
\paragraph{Definition 6.} For any tuple $w=(w_1,\ldots,w_t)$ and a subset $S\subseteq[t]$, we shall refer to the sum $\sum_{i\in S}w_i$ by $w_S$. And by $w|_S$, we will refer to the tuple obtained by considering only the elements of $w$ that are indexed by $S$. We denote by $\mathbb F_{\rm sm}[T]$ the set of set-multilinear polynomials over the tuple of sets of variables $T$.
The following is a simple result that establishes various useful properties of the relative rank measure.
\paragraph{Claim 7 ([23]).}
\begin{enumerate}
\item (Imbalance) Say $f\in\mathbb F_{\rm sm}[X(w)]$. Then, $\operatorname{relrk}_w(f)\leq2^{-|w_{[d]}|/2}$.
\item (Sub-additivity) If $f,g\in\mathbb F_{\rm sm}[X(w)]$, then $\operatorname{relrk}_w(f+g)\leq\operatorname{relrk}_w(f)+\operatorname{relrk}_w(g)$.
\item (Multiplicativity) Say $f=f_1f_2\cdots f_t$ and assume that for each $i\in[t]$, $f_i\in\mathbb F_{\rm sm}[X(w|_{S_i})]$, where $(S_1,\ldots,S_t)$ is a partition of $[d]$. Then
\[\operatorname{relrk}_w(f)=\prod_{i\in[t]}\operatorname{relrk}_{w|_{S_i}}(f_i).\]
\end{enumerate}
\subsection*{2.2 Inner Product Gadget}
We crucially need the following observation to construct the hard polynomials in Theorems 1 and 3.
\paragraph{Observation 8.} Let $n=2^k$ and $X_1=\{x_{1,1},\ldots,x_{1,n}\}$ and $X_2=\{x_{2,1},\ldots,x_{2,n}\}$ be two disjoint sets of variables. Then, for any symmetric word $w\in\{k,-k\}^2$ (i.e., where $w_1+w_2=0$) and for the inner product ``gadget'' $f=X_1\cdot X_2=\sum_{i=1}^n x_{1,i}x_{2,i}$, $\operatorname{relrk}_w(f)=1$ i.e., $M_w(f)$ is full-rank.
\section*{3 A Hard Set-multilinear Polynomial in VP}
\subsection*{3.1 Description of the Polynomial}
Let $d$ be an even integer and let $X=(X_1,\ldots,X_d)$ be a collection of sets of variables where each $|X_i|=n$, and similarly, let $Y=(Y_1,\ldots,Y_d)$ be a distinct collection of sets of variables where each $|Y_i|=n$. We shall refer to the $Y$-variables as the auxiliary variables. For $i$ and $j\in\{1,\ldots,d\}$, let $X_i\cdot X_j$ denote the inner-product quadratic form $\sum_{k=1}^n x_{ik}x_{jk}$. Here, we shall assume that $X_i=\{x_{i,1},\ldots,x_{i,n}\}$ and $Y_i=\{y_{i,1},\ldots,y_{i,n}\}$.
For two integers $i\in\mathbb N$ and $j\in\mathbb N$, we denote $[i,j]=\{k\in\mathbb N:i\leq k\text{ and }k\leq j\}$ and call such a set an interval. For every interval $[i,j]\subseteq[d]$, we define a polynomial $f_{i,j}(X,Y)\in\mathbb F_{\rm sm}[X_i,\ldots,X_j,Y_i,\ldots,Y_j]$ as follows:
\[f_{i,j}=\begin{cases}
y_{i,j}y_{j,i}(X_i\cdot X_j)&\text{if }j=i+1,\\
0&\text{if }j-i\text{ is even},\\
y_{i,j}y_{j,i}(X_i\cdot X_j)\cdot f_{i+1,j-1}+\displaystyle\sum_{r=i+1}^{j-1} f_{i,r}f_{r+1,j}&\text{otherwise}.
\end{cases}\]
\paragraph{Remark 9.} As described in Section 1.5, other than the use of the inner product gadget, one key difference between $f_{i,j}$ and the construction in [29] is that it uses fewer auxiliary variables. More specifically, while [29] had a ``fresh'' auxiliary variable for every choice of $i,r,j$ in the sum, we are unable to afford that not only because it destroys the set-multilinearity of the polynomial but most importantly, because of the aforementioned degree blow-up. This is also the reason why more straightforward attempts to ``set-multilinearize'' [29] such as by adding two ``copies'' $y_0$ and $y_1$ for each of their auxiliary variables $y$ (where intuitively $y_0$ and $y_1$ correspond to setting $y$ as $0$ or $1$ respectively in their argument) do not work.
The following is a more precise and general version of Theorem 3 that is stated in Section 1. We also incorporate Remark 4 here and show our lower bound for any degree $d\leq n$. Theorem 3 follows from the special case $d=n$.
\paragraph{Theorem 10.} Let $n=2^k$, and suppose $d\leq n$ be an even integer that is large enough\footnote{We only need $d$ to be larger than some absolute constant.}, and $1\leq\Delta\leq\log d/\log\log d$ be any positive integer. Let $X_i,Y_i$ denote the sets of $n$ variables $\{x_{i,j}:j\in[n]\}$ and $\{y_{i,j}:j\in[n]\}$ respectively and let $X,Y$ be the tuples $(X_1,\ldots,X_d)$ and $(Y_1,\ldots,Y_d)$. Then,
\begin{itemize}
\item there is a $\operatorname{poly}(n,d)$-size set-multilinear circuit computing $F_{n,d}=f_{1,d}(X,Y)$ as defined above,
\item any set-multilinear formula of product-depth $\Delta$ computing $F_{n,d}$ must have size at least $d^{\Omega(d^{1/\Delta}/\Delta)}$, and
\item further, any set-multilinear formula of arbitrary product-depth computing $F_{n,d}$ must have size at least $d^{\Omega(\log d)}$.
\end{itemize}
\subsection*{3.2 Proof of Hardness}
Note that the first item in Theorem 10 follows immediately from the recursive definition of $f_{1,d}$ (notice that there are only up to $d^2$ many distinct intervals of $[d]$). For proving the next two items, we invoke the symmetric word framework of [22]. The following couple of lemmas help establish that the relative rank measure with respect to symmetric words is (suitably) small for low-depth and general-depth set-multilinear formulas, respectively.
\paragraph{Lemma 11 ([22]).} Let $C$ be a set-multilinear formula of product-depth $1\leq\Delta\leq\log d/\log\log d$ of size at most $s$ which computes a polynomial (over any fixed field) that is set-multilinear with respect to the partition $(X_1,\ldots,X_d)$ where each $|X_i|=n$. Let $w\in\{k,-k\}^d$ be chosen uniformly at random. Then, we have
\[\operatorname{relrk}_w(C)\leq s\cdot2^{-kd^{1/\Delta}/20}\]
with probability at least $1-s\cdot d^{-d^{1/\Delta}/(12\Delta)}$.
\paragraph{Lemma 12 ([22]).} Let $F$ be a set-multilinear formula of size at most $s$ which computes a polynomial (over any fixed field) that is set-multilinear with respect to the partition $(X_1,\ldots,X_d)$ where each $|X_i|=n$. Let $w\in\{k,-k\}^d$ be chosen uniformly at random. Then, we have
\[\operatorname{relrk}_w(F)\leq s\cdot2^{-k\log d/20}\]
with probability at least $1-s\cdot d^{-\log d/60}$.
Next, we shall show that the hard polynomial $F_{n,d}$ in Theorem 10 has high relative rank (in fact, the maximum possible value -- $1$) with respect to a symmetric word. For this, we consider an alternate view of these polynomials and require the following notion. For an even integer $d$, define $\operatorname{Dyck}(d)$ to be the collection of all strings (called Dyck words) of length $d$ over symbols `(' and `)' that are well-matched in the natural way. More precisely, it is the collection of all strings $u$ of length $d$ such that all prefixes of $u$ contain no more $)$'s than $($'s and the total number of $($'s in $u$ equals the total number of $)$. For example, ``()()'' and ``(())'' belong to $\operatorname{Dyck}(4)$ but not ``(()(''. Note that for any `(' appearing in a Dyck word, there is a unique `)' which ``closes'' it. Given a Dyck word $u\in\operatorname{Dyck}(d)$, we call $(i,j)$ a matching parenthesis pair of $u$ if there is `(' in the $i$-th position of $u$ that is closed by a `)' in the $j$-th position of $u$ (clearly then, $j-i>0$ must be odd).
Given a string $u\in\operatorname{Dyck}(d)$ and the setup above for defining the polynomials $f_{i,j}$, we can associate to $u$ a product of inner products polynomial $\operatorname{IP}_u\in\mathbb F_{\rm sm}[X_i,\ldots,X_j]$ in the natural way: define $\operatorname{IP}_u$ to be the product of all $X_i\cdot X_j$ where $(i,j)$ is a matching parenthesis pair of $u$. For example, the strings ``()()'' and ``(())'' would correspond to the polynomials $(X_1\cdot X_2)\cdot(X_3\cdot X_4)$ and $(X_1\cdot X_4)\cdot(X_2\cdot X_3)$ respectively. We define $y_u$ analogously: it is the product of all $y_{i,j}\cdot y_{j,i}$ where $(i,j)$ is a matching parenthesis pair of $u$. So, if $u=\text{``(())''}\in\operatorname{Dyck}(4)$, then $y_u=y_{1,4}y_{4,1}y_{2,3}y_{3,2}$. The following observation then follows immediately from the recursive definition of $f_{i,j}$.
\paragraph{Observation 13.} For every interval $[i,j]\subseteq[d]$ where $j-i$ is odd, there exist constants $c_u\in\mathbb F$ corresponding to every $u\in\operatorname{Dyck}(j-i+1)$ such that
\[f_{i,j}(X,Y)=\sum_{u\in\operatorname{Dyck}(j-i+1)}c_u y_u\operatorname{IP}_u.\]\footnote{Strictly speaking, the indices within $\operatorname{IP}_u$ and $y_u$ here need to be ``translated'' appropriately to suit the interval $[i,j]$ (which may not necessarily be $[1,j-i+1]$).}
Moreover, if $\mathbb F$ has characteristic zero, then every $c_u\neq0$.
\paragraph{Claim 14.} Let $d$ be a positive even integer. For any $w\in\{-k,k\}^d$ with $w_{[d]}=0$ (i.e., $|P_w|=|N_w|$), there exists a Dyck work $u\in\operatorname{Dyck}(d)$ such that for every matching parenthesis pair $(i,j)$ of $u$, either $i\in P_w$ and $j\in N_w$, or $i\in N_w$ and $j\in P_w$.
\paragraph{Proof.} We prove this by induction on $d$. The base case $d=2$ is trivial as there is only a single matching parenthesis pair $(1,2)$ for which the given condition must indeed hold. Now, suppose $d>2$ and $w\in\{-k,k\}^d$ is a given word with $w_{[d]}=0$. Let us refer to $P_w$ and $N_w$ as the two ``parts'' of $w$ and take cases on the membership of $1$ and $d$ in these sets.
\begin{description}
\item[Case 1:] $1$ and $d$ are in different parts. By the induction hypothesis, there is a Dyck word $u'\in\operatorname{Dyck}(d-2)$ (corresponding to the subset $\{2,\ldots,d-1\}$ of indices) for which the desired condition holds. Hence, we can simply define $u$ to be the string ``$(u')$'' and the claim follows.
\item[Case 2:] $1$ and $d$ are in the same part. Notice that there exists an index $r$ such that both $w_{[1,r]}$ and $w_{[r+1,d]}$ are $0$. Then, we can define $u$ to be the concatenation of the two Dyck words that the induction hypothesis yields for the intervals $[1,r]$ and $[r+1,d]$ respectively and the claim follows.
\end{description}\hfill$\triangleleft$
\paragraph{Lemma 15.} Let $n=2^k$ and $d\leq n$ be an even integer. Over any field $\mathbb F$ of characteristic zero, the polynomial $F_{n,d}=f_{1,d}\in\mathbb F_{\rm sm}[X,Y]$ as defined above satisfies the following: For any $w\in\{-k,k\}^d$ with $w_{[d]}=0$, $M_w(F_{n,d})$ is full-rank when viewed as a matrix over the field $\mathbb F(Y)$, the field of rational functions over the $Y$ variables.
\paragraph{Proof.} Fix a word $w\in\{-k,k\}^d$ with $w_{[d]}=0$ and let $s\in\operatorname{Dyck}(d)$ be a Dyck word as given by Claim 14. By Observation 13, we know that $F=f_{1,d}$ has the form $\sum_{u\in\operatorname{Dyck}(d)}c_u y_u\operatorname{IP}_u$. Consider the polynomial $f$ obtained by plugging in $y_{i,j}=y_{j,i}=0$ for every $i,j$ such that $(i,j)$ is not a matching parenthesis pair of $s$, and $y_{i,j}=y_{j,i}=1$ for every $i,j$ such that $(i,j)$ is a matching parenthesis pair of $s$. Observe that the only surviving term from $F$ in $f$ is the one indexed by $s$. Therefore, to argue that $M_w(F)$ is full-rank over $\mathbb F(Y)$\footnote{We need to show that its determinant -- a polynomial in $\mathbb F[Y]$ -- is non-zero.}, it suffices to show that $M_w(c_s\operatorname{IP}_s)$ is full-rank. As $c_s\neq0$ by Observation 13, this follows from Observation 8 (the matrix of the inner product gadget has full rank), Claim 14 ($w$ ``splits'' every matching parenthesis pair of $s$), and Claim 7 (the multiplicativity of $\operatorname{relrk}_w$).\hfill$\triangleleft$
Let us return to the proof of the last two items of Theorem 10. Let $C$ be a set-multilinear formula of product depth $\Delta$ of size $s$ computing $F_{n,d}(X)$ (now interpreted as a formula over the field $\mathbb F(Y)$). Suppose $s<d^{d^{1/\Delta}/(24\Delta)}$. Then, by Lemma 11, with probability at least $1-d^{-d^{1/\Delta}/(24\Delta)}$,
\[\operatorname{relrk}_w(C)\leq s\cdot2^{-kd^{1/\Delta}/20}.\]
But now, we can condition on the event that $w_{[d]}=0$ (which occurs with probability $\Theta(1/\sqrt d)$) to establish the existence of a word $w\in\{-k,k\}^d$ with $w_{[d]}=0$ such that $w$ satisfies $\operatorname{relrk}_w(C)\leq s\cdot2^{-kd^{1/\Delta}/20}$. This is because of the asymptotic bound $1/\sqrt d\gg d^{-d^{1/\Delta}/(24\Delta)}$, which follows from the given constraints on the parameters $d,\Delta$. Therefore, by Lemma 15,
\[s\geq2^{kd^{1/\Delta}/20}\cdot\operatorname{relrk}_w(C)=n^{d^{1/\Delta}/20}\]
which contradicts the assumption that $s<d^{d^{1/\Delta}/(24\Delta)}$. Thus, we conclude that $s\geq d^{d^{1/\Delta}/(24\Delta)}=d^{\Omega(d^{1/\Delta}/\Delta)}$.
Similarly, to see the final item of Theorem 10, let $F$ be a set-multilinear formula of size $s$ computing $F_{n,d}$ (now interpreted as a formula over the field $\mathbb F(Y)$). Suppose $s<d^{\log d/120}$. Then, by Lemma 12, with probability at least $1-d^{-\log d/120}$,
\[\operatorname{relrk}_w(F)\leq s\cdot2^{-k\log d/20}.\]
But now, we can condition on the event that $w_{[d]}=0$ (which occurs with probability $\Theta(1/\sqrt d)$) to establish the existence of a word $w\in\{-k,k\}^d$ with $w_{[d]}=0$ such that $w$ satisfies $\operatorname{relrk}_w(F)\leq s\cdot2^{-k\log d/20}$. This is because of the trivial asymptotic bound $1/\sqrt d\gg d^{-\log d/120}$. Therefore, again by Lemma 15,
\[s\geq2^{k\log d/20}\cdot\operatorname{relrk}_w(F)=n^{\log d/20}\]
which contradicts the assumption that $s<d^{\log d/120}$. Thus, we conclude that $s\geq d^{\log d/120}=d^{\Omega(\log d)}$.
